\chapter{Conjunctivitis}
\index{Conjunctivitis}
\label{sec:Conjunctivitis}

\section{Overview}
Viral conjunctivitis, most commonly caused by adenovirus, results from viral infection of the conjunctiva.
\section{Causes}
Bacterial conjunctivitis results from bacterial infection (S. aureus, S. pneumoniae, H. influenzae, N. gonorrhoeae, C. trachomatis).     
\section{Risk Factors}

\begin{itemize}[noitemsep]
  \item Genetic predisposition and family history
  \item Advancing age
  \item Lifestyle factors (diet, exercise, smoking, alcohol)
  \item Environmental exposures
  \item Underlying medical conditions
  \item Medication use
\end{itemize}
\section{Signs and Symptoms}

\subsection{Common symptoms}
\begin{itemize}[noitemsep]
  \item Symptoms of viral conjunctivitis include acute onset, watery discharge, and preauricular lymphadenopathy
  \item Bacterial conjunctivitis presents with purulent discharge and eyelids matted shut
  \item Allergic conjunctivitis features bilateral itching, tearing, and stringy mucus.
  \item Pain or discomfort in the affected area
  \end{itemize}

\subsection{Emergency warning signs}
\begin{itemize}[noitemsep]
  \item Severe or worsening symptoms of conjunctivitis require immediate medical evaluation
\end{itemize}
\section{Progression and Stages}
Conjunctivitis is an inflammation of the conjunctiva and the most common eye disease worldwide, classified as infectious (viral, bacterial), allergic, or chemical or irritant. The condition may progress through different stages of severity over time. Early stages often involve mild or subtle symptoms that may be overlooked. Moderate stages involve more noticeable symptoms affecting daily function. Advanced stages may involve significant impairment and complications.
\section{Complications}
\begin{itemize}[noitemsep]
  \item Disease progression leading to worsening symptoms
  \item Secondary infections or injuries
  \item Side effects of treatment
  \item Reduced quality of life
  \item Emotional and psychological impact
\end{itemize}
\section{Diagnosis}
Doctors usually diagnose this based on your symptoms and a physical exam, with Gram stain and culture used for hyperacute or bacterial cases. Comprehensive medical history and physical examination Laboratory tests to assess organ function and rule out other conditions Imaging studies to visualize affected structures Specialized testing depending on the affected system Consultation with relevant specialists Monitoring of disease progression over time

\begin{itemize}[noitemsep]
  \item Comprehensive medical history and physical examination
  \item Laboratory tests to assess organ function
  \item Imaging studies to visualize affected structures
  \item Specialized testing depending on the affected system
  \item Consultation with relevant specialists
\end{itemize}
\section{Prevention}
\begin{itemize}[noitemsep]
  \item Maintain a healthy lifestyle with balanced diet and exercise
  \item Avoid known risk factors
  \item Attend regular medical check-ups for early detection
  \item Follow recommended screening guidelines
\end{itemize}
\section{Treatment Options}
\begin{itemize}[noitemsep]
  \item Treatment typically involves etiologic: supportive for viral, topical antibiotics for bacterial, and topical antihistamine or mast cell stabilizers for allergic
  \item Gonococcal conjunctivitis requires intramuscular ceftriaxone.
  \item Medications to manage symptoms and address underlying mechanisms
  \item Lifestyle modifications and self-care measures
  \item Physical therapy and rehabilitation
  \item Surgical intervention when indicated
  \item Regular monitoring and follow-up care
\end{itemize}
\section{Living with It}
\begin{itemize}[noitemsep]
  \item Follow your treatment plan as prescribed
  \item Maintain regular follow-up with your healthcare provider
  \item Make necessary lifestyle adjustments
  \item Seek support from family, friends, or support groups
  \item Stay informed about your condition
\end{itemize}
\section{Outlook}
\begin{itemize}[noitemsep]
  \item Most people recover well with appropriate treatment
  \item Untreated gonococcal conjunctivitis can lead to corneal perforation within 24 hours.
\end{itemize}
